\paragraph{Author descriptions and Kolmogorov complexity (Section 2.4).}
Throughout we assume a fixed pairing function $\langle\cdot\rangle$ and a standard encoding of Turing machines $M$ by binary strings, with $\widetilde M$ denoting the encoding of $M$.
We fix a universal Turing machine $U$ that for any Turing machine $M$ and input $x$ takes as input $\langle\widetilde M,x\rangle$ and outputs $M(x)$. The Kolmogorov complexity $K(q)$ of $q$ (with respect to $U$), denoted $K(q)$, is the minimum length of $y$ such that $U(y)=q$. We define the Kolmogorov complexity of a machine $M$ to be $K(\widetilde M)$.
\paragraph{Proposition 11 (Symmetry of information).}
Let $x,y$ be strings. Then $K(x;y)=K(x)+K(y\mid x)+O(\log(K(x,y)))$.
\paragraph{Author random-string promise problems (Section 2.5).}
In the $R_K$ problem, the instance is a string $x$, and the question is whether $K(x)\geq|x|/2$. Given a function $g:\mathbb N\to\mathbb N$, the $R_K$ problem with gap $g$ is the promise problem whose NO instances are strings $x$ with $K(x)\geq|x|/2$ and whose YES instances are strings $x$ with $K(x)<|x|/2-g(|x|)$. Given a function $\alpha:\mathbb N\to\mathbb R$ such that $\alpha(n)\geq1$ for all $n$, the $R_K$ problem with multiplicative gap $\alpha$ is the promise problem whose NO instances are strings $x$ with $K(x)\geq|x|/2$ and whose YES instances are strings $x$ such that $K(x)<|x|/(2\alpha(|x|))$.
